% ============================================================================
% 02_diagrama_casos_uso.tex — Cuerpo del entregable "Diagrama de Casos de Uso"
% ============================================================================
\section{Diagrama de Casos de Uso}

\subsection{Actores}

La Tabla~\ref{tab:actores-cu} identifica los actores del sistema LSA Salud,
siguiendo la notación UML~2.5. Todos son actores primarios: cada uno
persigue un objetivo propio al interactuar con el sistema. ACT-02 y ACT-03
son decisiones de modelado explícitas frente al Registro de Interesados
oficial: se justifican en detalle en \texttt{IDS.md}, Sección~1.

\begin{table}[H]
    \centering
    \caption{Actores de LSA Salud}
    \label{tab:actores-cu}
    \begin{tabularx}{\textwidth}{l l X}
        \toprule
        \textbf{ID} & \textbf{Actor} & \textbf{Descripción} \\
        \midrule
        ACT-01 & Paciente & Persona sorda o hipoacúsica que se comunica mediante LSA frente a la cámara de la Estación del Paciente. Corresponde al interesado oficial ``Personas que se comunican con LSA''. \\
        ACT-02 & Personal de Recepción o Enfermería & Personal de recepción o enfermería que lee la traducción del paciente y responde por texto o por voz. Generaliza a los interesados oficiales ``Recepcionista'' y ``Enfermero/a'', que comparten requisito, interés y estrategia de gestión en el Registro de Interesados. \\
        ACT-03 & Desarrollador de IA / Equipo de Proyecto & Personal técnico que amplía el vocabulario de señas y reentrena el modelo de reconocimiento. Decisión de diseño del equipo: el Registro de Interesados no define un rol de usuario final equivalente; el trabajo de dataset y entrenamiento figura en el EDT como tarea del equipo del proyecto. \\
        \bottomrule
    \end{tabularx}
\end{table}

No se modelan actores de sistema secundarios (por ejemplo, el modelo de
lenguaje local o el motor de síntesis de voz). Aunque LSA Salud depende de
ellos, ninguno cumple la definición de actor de Larman/UML: no tienen
iniciativa propia ni interactúan directamente con una persona, sino que son
invocados de forma sincrónica por el propio backend como parte de la
implementación de los casos de uso. Se documentan como dependencias
tecnológicas en la Especificación de Requisitos de Software (secciones 2.4
y 3.4), no como actores del diagrama. Tampoco se modela un actor para el
consumo de las métricas de auditoría de CU-11 (RF-04 oficial): la
justificación completa está en \texttt{IDS.md}, Sección~1.

\subsection{Diagrama}

El diagrama sigue el formato de referencia de la cátedra: un único
rectángulo de límite del sistema titulado ``LSA Salud'', sin paquetes
internos, con los once casos de uso como elipses de fondo amarillo pálido y
borde fino, y los tres actores como figuras de palito fuera del límite,
unidos por líneas rectas. Los casos de uso se ordenan en tres columnas: a la
izquierda, los que inicia el Paciente (CU-09, CU-01 y sus inclusiones CU-02
y CU-03); en el centro, los compartidos por ambas estaciones (CU-10, que
extiende a CU-01 y CU-04, y CU-11, incluido por ambos); a la derecha, los del
Personal de Recepción o Enfermería (CU-04 y su extensión CU-05) y los del
Desarrollador de IA / Equipo de Proyecto (CU-06 a CU-08). El Paciente se
ubica a la izquierda del límite y los demás actores a la derecha.

\diagrama{02_casos_uso}{\textwidth}{Diagrama de Casos de Uso de LSA Salud}{fig:casos-uso}

\subsection{Relaciones modeladas}

\begin{itemize}
    \item \textbf{CU-01 «include» CU-02 y CU-03:} comunicarse mediante LSA
    siempre implica traducir la secuencia de señas a una oración y, si la
    síntesis de voz está disponible, sintetizarla en audio. Ambas son
    subrutinas obligatorias del flujo principal de CU-01, no
    comportamiento opcional, por lo que se modelan con \texttt{include}.
    \item \textbf{CU-09 es precondición de CU-01:} el paciente debe otorgar
    el consentimiento explícito de cámara (RF-03 oficial, diseñado y
    pendiente de implementación) antes de que exista cualquier sesión de
    reconocimiento. Se modela como precondición y no como \texttt{include}
    porque ocurre una sola vez por sesión, no en cada ejecución de CU-01.
    \item \textbf{CU-05 «extend» CU-04:} dictar la respuesta por voz es una
    forma \emph{alternativa y opcional} de completar CU-04 (Responder al
    Paciente). El personal puede optar por escribir su respuesta (flujo
    principal de CU-04) o activar el dictado por voz, que extiende ese flujo
    en el punto de extensión ``ingreso de texto''. Por tratarse de una
    variante opcional del comportamiento base, corresponde \texttt{extend}
    y no \texttt{include}.
    \item \textbf{CU-10 «extend» CU-01 y CU-04:} reportar una traducción
    incorrecta o ambigua (RF-06 oficial, diseñado y pendiente de
    implementación) es opcional y puede dispararse tanto sobre un mensaje
    del paciente (dentro de CU-01) como sobre uno del personal (dentro de
    CU-04), en el punto de extensión ``revisión de mensaje'' de cada uno.
    \item \textbf{CU-01 y CU-04 «include» CU-11:} registrar el evento de
    auditoría anónimo (RF-04 oficial, diseñado y pendiente de
    implementación) es una subrutina obligatoria de ambos casos de uso, sin
    comportamiento opcional ni actor humano que la dispare --- de ahí que se
    modele como \texttt{include} y no como \texttt{extend}.
    \item \textbf{CU-06 «include» CU-07:} capturar una muestra de
    entrenamiento siempre requiere seleccionar o crear la etiqueta de la
    seña que se va a grabar, por lo que la gestión del vocabulario es una
    subrutina obligatoria de la captura.
    \item \textbf{CU-08 (Entrenar el Modelo) es independiente} de CU-06 en el
    diagrama: el reentrenamiento del modelo LSTM se ejecuta de forma
    manual y offline (\texttt{train\_lstm.py}) sobre el conjunto acumulado
    de muestras, no como una invocación síncrona disparada por cada
    captura individual. La relación entre ambos se documenta como
    precondición en el Modelo de Casos de Uso (Entregable de Modelos de
    Casos de Uso), no como flecha de inclusión.
\end{itemize}

No se modelan generalizaciones de actores: en el estado actual del sistema
no existe una jerarquía real entre Paciente, Personal de Recepción o
Enfermería y Desarrollador de IA / Equipo de Proyecto —cada uno tiene
objetivos, precondiciones y estaciones de trabajo distintos—, por lo que
introducir un actor abstracto común no aportaría claridad. CU-11 no tiene
actor principal: es una subfunción interna sin interacción humana directa
(ver \texttt{IDS.md}, Sección~1).
